% ----------------------------------------------------
% Projeto: UnDF-TeX
% Arquivo: 02-referencial-teorico.tex
%
% Copyright (c) 2026
% Universidade do Distrito Federal (UnDF)
%
% Este arquivo faz parte do projeto UnDF-TeX.
%
% Licenciado sob a LaTeX Project Public License (LPPL),
% versão 1.3c ou posterior.
%
% Implementação técnica baseada nas normas:
% ABNT NBR 14724:2024
% ABNT NBR 10520:2023
%
% Este projeto não reproduz o texto integral
% das normas da ABNT.
% ----------------------------------------------------

\chapter{Referencial teórico}

A produção acadêmica exige rigor na construção das ideias e clareza na forma de apresentá-las. Segundo \textcite{severino2017}, o método não se separa da inteligibilidade do texto, pois organizar bem o trabalho é também organizar o pensamento.\footnote{Esta nota de rodapé exemplifica um recurso textual permitido quando necessário ao esclarecimento pontual do texto.}

A citação direta curta pode ser exemplificada assim: \citedir[p. 293]{derrida1967}{a desconstrução do logocentrismo não é uma psicanálise da filosofia [...]}. Já a citação de citação pode ser exemplificada assim:
\apudbibonline[p. 13][p. 25]{freire1994}{streck2017}.

A citação direta longa, por sua vez, deve aparecer destacada, sem aspas, com fonte menor, espaçamento simples e recuo em relação à margem esquerda. O bloco a seguir usa um texto didático do próprio projeto apenas para demonstrar a forma:

\begin{citacaolonga}[p. 1]{undftex2026}
Transcreva aqui, literalmente, um trecho com mais de três linhas. Este bloco demonstra a forma de apresentação da citação direta longa: fonte menor, espaçamento simples, recuo à esquerda e ausência de aspas. Substitua este texto por uma citação real, conferida na fonte consultada.
\end{citacaolonga}
